% Figure — Du workflow d'évaluation au calcul des métriques (phase 1).
% Prolonge la figure « Workflow d'évaluation n8n » : mêmes couleurs, mêmes formes.
% Aucune bibliothèque TikZ supplémentaire requise (graphicx pour \resizebox).
\providecolor{couleurCorpus}{RGB}{241,247,244}
\providecolor{couleurCorpusBord}{RGB}{26,122,76}
\providecolor{couleurBleu}{RGB}{31,95,168}
\providecolor{couleurBleuFond}{RGB}{242,245,250}
\providecolor{couleurOrange}{RGB}{252,229,205}
\providecolor{couleurOrangeBord}{RGB}{214,155,80}
\providecolor{couleurN8n}{RGB}{234,75,113}
\resizebox{\linewidth}{!}{%
\begin{tikzpicture}[
  x=1cm, y=1cm, >=latex,
  f/.style={->, very thick, draw=couleurBleu},
  leg/.style={font=\small, align=center, anchor=north}
]
% ---------------- rappel du workflow n8n ----------------
\node[draw=couleurN8n, very thick, rounded corners=8pt, fill=white,
      align=center, text width=2.7cm, minimum height=2.2cm, inner sep=4pt,
      font=\small] (n8n) at (-4.3,-1.6)
  {\textbf{Workflow n8n}\\(figure précédente)\\126 exécutions\\par modèle};

% ---------------- entrée 1 : fichiers de résultats JSONL ----------------
\begin{scope}[shift={(-0.4,-0.3)}]
  \foreach \d in {0.30, 0.15, 0}{
    \filldraw[draw=couleurOrangeBord, fill=couleurOrange, thick]
      (\d,-1.15+\d) -- (\d,1.15+\d) -- (1.70+\d,1.15+\d) -- (2.35+\d,0.50+\d) -- (2.35+\d,-1.15+\d) -- cycle;
    \filldraw[draw=couleurOrangeBord, fill=white, thick]
      (1.70+\d,1.15+\d) -- (1.70+\d,0.50+\d) -- (2.35+\d,0.50+\d);
  }
  \node[font=\bfseries\small, text=couleurOrangeBord] at (0.85,0.85) {JSONL};
  \node[font=\scriptsize\ttfamily, align=left, anchor=north west] at (0.12,0.45)
    {\{"meta"…\},\\\ "mapping"…,\\\ "usage"…\}};
\end{scope}
\node[leg] at (0.95,-1.65)
  {\textbf{Fichiers de résultats}\\\texttt{resultats\_\_<modèle>.jsonl}};

% ---------------- entrée 2 : vérité de terrain ----------------
\begin{scope}[shift={(-0.25,-4.6)}]
  \filldraw[draw=couleurCorpusBord, fill=couleurCorpus, thick]
    (0,-1.15) -- (0,1.15) -- (1.70,1.15) -- (2.35,0.50) -- (2.35,-1.15) -- cycle;
  \filldraw[draw=couleurCorpusBord, fill=white, thick]
    (1.70,1.15) -- (1.70,0.50) -- (2.35,0.50);
  \node[font=\scriptsize\ttfamily, align=left, anchor=north west] at (0.12,0.45)
    {S000 : TA0102,\\\ \ T0846\\S001 : …};
\end{scope}
\node[leg] at (0.90,-5.95)
  {\textbf{Vérité de terrain}\\identifiant attendu\\aux trois niveaux};

% ---------------- le script ----------------
\node[draw=couleurBleu, very thick, rounded corners=8pt, fill=couleurBleuFond,
      align=center, text width=5.6cm, minimum height=4.4cm, inner sep=6pt,
      font=\small] (script) at (7.3,-2.2)
  {\textbf{\texttt{calculer\_metriques.py}}\\[5pt]
   comparaison de chaque\\prédiction à l'attendu\\[3pt]
   $\downarrow$\\[3pt]
   \textbf{étiquetage} : VP, FP, H, O, A\\[3pt]
   $\downarrow$\\[3pt]
   \textbf{métriques} : précision,\\rappel et F1 aux trois niveaux,\\dans les deux bases};

% ---------------- flèches d'entrée ----------------
\draw[->, very thick, draw=couleurN8n] (n8n.east) -- (-0.55,-0.9);
\draw[->, very thick, draw=couleurOrangeBord] (2.30,-0.30) -- (4.25,-1.30);
\draw[->, very thick, draw=couleurCorpusBord] (2.30,-4.60) -- (4.25,-3.10);

% ---------------- sortie 1 : rapport texte ----------------
\begin{scope}[shift={(12.1,0.1)}]
  \filldraw[draw=couleurBleu, fill=couleurBleuFond, thick]
    (0,-1.15) -- (0,1.15) -- (1.70,1.15) -- (2.35,0.50) -- (2.35,-1.15) -- cycle;
  \filldraw[draw=couleurBleu, fill=white, thick]
    (1.70,1.15) -- (1.70,0.50) -- (2.35,0.50);
  \node[font=\bfseries\small, text=couleurBleu] at (0.80,0.85) {TXT};
  \foreach \y in {0.20, -0.10, -0.40, -0.70}{
    \draw[draw=couleurBleu, thick] (0.25,\y) -- (2.05,\y);
  }
\end{scope}
\node[leg] at (13.3,-1.25)
  {\textbf{Rapport texte}\\lecture directe,\\modèle par modèle};

% ---------------- sortie 2 : fichier CSV ----------------
\begin{scope}[shift={(12.1,-4.5)}]
  \filldraw[draw=couleurCorpusBord, fill=couleurCorpus, thick]
    (0,-1.15) -- (0,1.15) -- (1.70,1.15) -- (2.35,0.50) -- (2.35,-1.15) -- cycle;
  \filldraw[draw=couleurCorpusBord, fill=white, thick]
    (1.70,1.15) -- (1.70,0.50) -- (2.35,0.50);
  \node[font=\bfseries\small, text=couleurCorpusBord] at (0.80,0.85) {CSV};
  \draw[draw=couleurCorpusBord, thick] (0.25,0.25) rectangle (2.05,-0.85);
  \draw[draw=couleurCorpusBord, thick] (0.25,-0.10) -- (2.05,-0.10);
  \draw[draw=couleurCorpusBord, thick] (0.25,-0.48) -- (2.05,-0.48);
  \draw[draw=couleurCorpusBord, thick] (1.15,0.25) -- (1.15,-0.85);
\end{scope}
\node[leg] at (13.3,-5.85)
  {\textbf{Fichier CSV}\\tableaux de résultats\\repris dans ce chapitre};

% ---------------- flèches de sortie ----------------
\draw[f] (10.35,-1.30) -- (12.00,-0.30);
\draw[f] (10.35,-3.10) -- (12.00,-4.10);
\end{tikzpicture}%
}
